\section{Пример листинга в приложении.}\label{sec:appB}

\hiddensubsection{Раздел приложения}

\hiddensubsection{Раздел приложения}

\hiddensubsection{Раздел приложения}

\lipsum[1]

\begin{longlisting}
    \caption{Пример длинного листинга}
	\begin{minted}[autogobble,breaklines,extrakeywordstype=vector,extrakeywordsconstant={true,false}]{c++}
		#include <iostream>
		#include <vector>
		#include <cmath>
		using namespace std;
		
		vector<int> sieveOfEratosthenes(int n) {
			vector<bool> isPrime(n + 1, true);
			vector<int> primes;
			
			isPrime[0] = isPrime[1] = false;
			
			for (int i = 2; i * i <= n; i++) {
				if (isPrime[i]) {
					for (int j = i * i; j <= n; j += i) {
						isPrime[j] = false;
					}
				}
			}
			
			for (int i = 2; i <= n; i++) {
				if (isPrime[i]) {
					primes.push_back(i);
				}
			}
			
			return primes;
		}
		
		int main() {
			int limit;
			cout << "Введите верхнюю границу для поиска простых чисел: ";
			cin >> limit;
			
			vector<int> primes = sieveOfEratosthenes(limit);
			
			cout << "Найдено простых чисел: " << primes.size() << endl;
			cout << "Все простые числа до " << limit << ":" << endl;
			
			for (int i = 0; i < primes.size(); i++) {
				cout << primes[i] << " ";
				if ((i + 1) % 10 == 0) cout << endl; // 10 чисел в строке
			}
			cout << endl;
			
			return 0;
		}	
	\end{minted}

\end{longlisting}

